\uuid{Vs0A}
\exo7id{3845}
\titre{exo7 3845}
\auteur{quercia}
\organisation{exo7}
\datecreate{2010-03-11}
\isIndication{false}
\isCorrection{false}
\chapitre{Continuité, limite et étude de fonctions réelles}
\sousChapitre{Continuité : théorie}
\cpge{mpsi-mp2i.fonctions-valeurs-reelles-complexes.a, pcsi.techniques-calcul-fonctions.a, ptsi.techniques-calcul-fonctions.a, mpsi-mp2i.limites-continuite.a, pcsi.limites-continuite.a, ptsi.limites-continuite.a, mpsi-mp2i.limites-continuite.c, pcsi.limites-continuite.c, ptsi.limites-continuite.c, bcpst1.fonctions-reelles-usuelles.a, bcpst1.limites-continuite.a, bcpst1.limites-continuite.c}
\module{Analyse}
\niveau{L1}
\difficulte{}

\contenu{
\texte{
Soit $f : \R \to \R$ et $T > 0$. On suppose que $f$ est $T$-périodique cad :
$\forall\ x \in \R,\ f(x+T) = f(x)$.
}
\begin{enumerate}
    \item \question{Si $f$ possède une limite en $+\infty$, montrer que $f$ est constante.}
    \item \question{Si $f$ est continue non constante, montrer que $f$ a une plus petite période.}
    \item \question{Si $f$ est continue, montrer que $f$ est bornée et atteint ses bornes.}
\end{enumerate}
}
